\documentclass{article}
\PassOptionsToPackage{numbers,sort,compress}{natbib}
\usepackage[main,final]{neurips_2025}
\usepackage[T2A]{fontenc}
\usepackage[utf8]{inputenc}
\usepackage[russian]{babel}
\usepackage{amsmath,amssymb}
\usepackage[unicode,hidelinks]{hyperref}
\usepackage{graphicx}
\renewcommand{\rmdefault}{cmr}
\renewcommand{\sfdefault}{cmss}
\makeatletter
\renewcommand{\@noticestring}{Выпускная квалификационная работа.}
\makeatother
\setlength{\emergencystretch}{2em}
\renewenvironment{abstract}
{\vskip 0.075in\centerline{\large\bfseries Аннотация}\vspace{0.5ex}\begin{quote}}
{\par\end{quote}\vskip 1ex}

\title{Влияние различий средних на обнаружение корреляционных паттернов при интерпретации кластеров}
\author{}
\date{}

\begin{document}
\maketitle

\begin{abstract}
Интерпретация кластеров требует выявления не только различий отдельных признаков, но и зависимостей между ними, которые могут не отражаться в рейтингах важности.
Цель работы состоит в оценке влияния межгрупповых различий средних на обнаружение признаков, участвующих в корреляционном паттерне, и в сравнении способов измерения их важности.
Для этого используется контролируемый синтетический эксперимент с двумя заранее заданными группами, в которых корреляции пары признаков различаются, а их средние и дисперсии совпадают; разность средних дополнительного признака последовательно увеличивается при неизменном распределении исходной пары.
Сравниваются случайный лес с перестановочной важностью по доле правильных ответов и логарифмической функции потерь, локальные интерпретируемые объяснения, метод главных компонент и прямое сравнение корреляций; основной показатель определяется как частота одновременного включения обоих признаков пары в тройку наиболее важных по результатам двадцати независимых повторов.
В базовом эксперименте при увеличении разности средних дополнительного признака от нуля до четырёх внутригрупповых стандартных отклонений эта частота у случайного леса снизилась со ста до нуля процентов для важности по доле правильных ответов и со ста до семидесяти процентов для важности по логарифмическим потерям.
Метод главных компонент и прямое сравнение корреляций сохраняли оба признака во всех повторах, тогда как выбранный вариант локальных объяснений редко выделял пару уже при отсутствии различий средних.
Полученные предварительные результаты показывают необходимость учитывать способ оценки важности при интерпретации структуры групп и отделять полноту выделения участников паттерна от качества объяснения предсказаний классификатора.
\end{abstract}

\section{Введение}

При интерпретации кластеров требуется установить, какие свойства данных характеризуют выделенные группы. Такие свойства включают различия средних отдельных признаков и различия зависимостей между ними. Цель исследования --- определить, как усиление первого типа различий влияет на обнаружение признаков, участвующих во втором.

Для интерпретации готового разбиения применяются модели, предсказывающие принадлежность к кластеру, локальные объяснения и методы снижения размерности. Connor и соавторы~\cite{connor2026} сравнивают случайный лес, LIME и метод главных компонент на данных с известными паттернами. LIME строит локальное приближение предсказаний модели~\cite{ribeiro2016}, а разложение DIP разделяет самостоятельные вклады признаков, взаимодействия и зависимости~\cite{konig2025}.

Эти подходы оценивают разные характеристики: влияние признака на прогноз, поведение модели в окрестности объекта или структуру вариации данных. Поэтому высокое качество классификации и корректное объяснение предсказания не гарантируют полного описания межгрупповых различий. Кроме того, стандартный случайный лес может недостаточно учитывать некоторые взаимодействия~\cite{blum2025}, что необходимо отделять от ограничений самой процедуры интерпретации.

В работе проводится сравнение при фиксированном корреляционном паттерне и изменяемом различии средних дополнительного признака. Такая постановка позволяет проследить изменение рейтингов при изменении одного фактора. Задачи исследования --- измерить частоту выделения участников паттерна, сопоставить способы оценки важности и проверить устойчивость результатов к параметрам данных и методов.

Первый эксперимент проводится на синтетических данных с заданными групповыми метками. Для каждого независимого повтора создаются обучающая и тестовая выборки, затем изменяется только сдвиг среднего дополнительного признака. Предварительный результат --- различная чувствительность двух перестановочных оценок важности одного классификатора к этому сдвигу. Практический смысл сравнения заключается в выборе способов интерпретации, сохраняющих информацию об исследуемой структуре групп.

\section{Постановка задачи}

Даны обучающая и независимая тестовая выборки с фиксированными групповыми метками. Требуется оценить, как изменение различий средних одного признака влияет на выделение другой пары признаков, корреляции которой различаются между группами.

Рассматриваются две равные группы A и B. Признаки X и Y имеют стандартные нормальные распределения в обеих группах; их корреляция равна 0,8 в группе A и 0 в группе B. Признак Z независим от X и Y внутри каждой группы и имеет единичную дисперсию. Средние Z сдвигаются симметрично относительно нуля, а разность между ними последовательно увеличивается. Распределение пары X--Y внутри групп при этом не меняется. Остальные признаки представляют собой независимый стандартный нормальный шум.

Каждый метод формирует рейтинг признаков. Основной критерий качества --- доля независимых повторов, в которых X и Y одновременно входят в тройку наиболее важных. Дополнительно оцениваются ранги признаков и качество классификатора на тестовой выборке. Метрика характеризует выделение участников паттерна, а не восстановление самой зависимости.

Важность случайного леса оценивается изменением доли правильных ответов или логарифмических потерь после перемешивания признака на тестовой выборке. Для LIME используются усреднённые абсолютные веса локальных объяснений, для метода главных компонент --- абсолютные нагрузки первой компоненты. Прямое сравнение корреляций служит базовым методом. На первом этапе группы заданы заранее, что позволяет исследовать интерпретацию отдельно от алгоритма кластеризации.

\section{Теоретическая основа и методы}

\subsection{Различия средних и корреляционные паттерны}

Разность средних характеризует смещение значений одного признака между группами. Корреляция Пирсона описывает направление и силу линейной связи двух признаков внутри группы: положительная корреляция соответствует их совместному увеличению, отрицательная --- изменению в противоположных направлениях. Нулевая корреляция означает отсутствие линейной связи, но в общем случае не гарантирует независимости.

В данной работе корреляционным паттерном называется различие корреляций одной пары признаков между группами. При этом распределения каждого из двух признаков по отдельности могут совпадать. Поэтому анализ средних не позволяет обнаружить такое различие, а рейтинг отдельных признаков не описывает саму зависимость~\cite{connor2026}.

\subsection{Случайный лес и перестановочная важность}

Случайный лес (Random Forest, RF) представляет собой ансамбль деревьев решений, обучаемых с использованием случайных подвыборок объектов и признаков~\cite{breiman2001}. В эксперименте лес обучается предсказывать готовые групповые метки. Классификация используется как вспомогательная задача: её объяснение рассматривается как один из способов описания различий групп.

Перестановочная важность (permutation importance) оценивает, насколько ухудшается качество обученной модели после случайного перемешивания значений одного признака между тестовыми объектами. Модель при этом не переобучается. Большое ухудшение указывает на существенную роль признака для выбранного критерия качества. Зависимости между признаками влияют на такую оценку, поскольку перемешивание нарушает их совместное распределение~\cite{strobl2008}.

Используются две меры качества одного и того же леса. Accuracy --- доля правильных ответов --- изменяется при изменении предсказанного класса. Log loss --- логарифмические потери --- учитывает также вероятность, присвоенную правильному классу. Например, уменьшение вероятности правильного класса с 0,9 до 0,6 не меняет accuracy при сохранении ответа, но увеличивает log loss. Следовательно, один признак может получить разные оценки важности в зависимости от выбранной меры.

\subsection{Локальные объяснения LIME}

LIME (Local Interpretable Model-agnostic Explanations) объясняет отдельное предсказание с помощью простой модели, приближающей поведение исходного классификатора в окрестности объекта~\cite{ribeiro2016}. Для этого создаются возмущённые варианты объекта, вычисляются ответы классификатора и подбирается взвешенное локальное приближение. Коэффициенты этого приближения служат оценками влияния признаков на конкретное предсказание.

В эксперименте LIME объясняет вероятности группы A, выдаваемые тем же случайным лесом. Для получения общего рейтинга усредняются абсолютные веса признаков по выбранным тестовым объектам. Такое усреднение является принятой процедурой эксперимента, а не гарантией глобального объяснения. Простое локальное приближение может недостаточно отражать совместное действие X и Y; поэтому отдельно оценивается его согласованность с ответами леса. Низкое качество приближения ограничивает интерпретацию полученного рейтинга.

\subsection{Метод главных компонент}

Метод главных компонент (Principal Component Analysis, PCA) находит линейные комбинации признаков, последовательно объясняющие наибольшую дисперсию данных~\cite{sklearnpca}. При использовании этого метода групповые метки не служат целевой переменной классификации. В работе рейтинг определяется абсолютными значениями коэффициентов первой компоненты после стандартизации признаков.

Рассматриваются общая стандартизация обучающей выборки и отдельная стандартизация внутри каждой группы. Во втором случае абсолютные коэффициенты усредняются по группам. Внутригрупповое центрирование устраняет добавленный сдвиг Z, поэтому этот вариант используется как контроль. Высокая нагрузка компоненты характеризует структуру вариации, но сама по себе не доказывает, что признак различает группы.

\subsection{Базовый метод и критерий сравнения}

Прямой базовый метод сравнивает корреляции каждой пары признаков между заданными группами. Разность оценивается после преобразования Фишера с учётом размеров групп; оценкой отдельного признака служит максимальное значение среди содержащих его пар. Такой метод непосредственно использует исследуемую характеристику и позволяет сопоставить модельные объяснения с анализом структуры данных.

Все методы сравниваются по попаданию обоих участников известного паттерна в top-3. Этот критерий оценивает полноту выделения признаков. Он не тождественен качеству объяснения предсказаний: Z является полезным предиктором групп, и его высокая важность корректна даже тогда, когда X и Y занимают более низкие позиции.

\section{Базовый эксперимент и первые результаты}

Использованы выборки по 600 объектов для обучения и независимой проверки, с равным числом объектов каждой группы. Размерность 10, число повторов 20, значения разности средних Z, обозначенной через $\Delta$: 0; 0,5; 1; 2; 3; 4. Внутри повтора используются одинаковые исходные объекты и случайные величины при всех значениях $\Delta$.

В таблице приведены значения основной метрики для крайних значений сдвига. Вариант метода главных компонент со стандартизацией внутри групп служит контролем: такая обработка устраняет добавленный сдвиг среднего.

\begin{table}[ht]
\centering
\small
\begin{tabular}{lrr}
\hline
Метод & $\Delta=0$ & $\Delta=4$ \\
\hline
Случайный лес: важность по accuracy & $100\%$ & $0\%$ \\
Случайный лес: важность по log loss & $100\%$ & $70\%$ \\
LIME: усреднённые абсолютные веса & $5\%$ & $5\%$ \\
Главные компоненты: общая стандартизация & $100\%$ & $100\%$ \\
Главные компоненты: внутри групп & $100\%$ & $100\%$ \\
Прямое сравнение корреляций & $100\%$ & $100\%$ \\
\hline
\end{tabular}
\caption{Включение обоих признаков $X,Y$ в тройку наиболее важных.}
\end{table}

Средняя доля правильных ответов случайного леса выросла с $57{,}8\%$ до $97{,}8\%$. При этом выделение пары по перестановочной важности существенно зависело от выбранной меры качества. Результат LIME нельзя объяснить только сдвигом $Z$, поскольку пара редко выделялась уже при $\Delta=0$.

Число повторов ограничивает точность оценок: для частот $0\%$, $70\%$ и $100\%$ интервалы Wilson уровня $95\%$ составляют соответственно $[0;16{,}1]\%$, $[48{,}1;85{,}5]\%$ и $[83{,}9;100]\%$. Полученные результаты относятся к выбранному генератору и настройкам методов. Далее требуется проверить другие силы корреляции, объёмы выборки и размерности данных.

\bibliographystyle{unsrtnat}
\bibliography{references}
\end{document}
